\section{機能群の連携、計算フローと数値実装}
\label{sec:implementation}
前節までに述べた物理式と推定法は、入力の準備、状態ごとの反復評価、最終結果の出力という三つの段階で実行される。本節では、それぞれの機能が何を受け取り、次の機能へ何を渡すかを整理する。続いて、積分、行列解法、観測格子、誤差処理を具体的に示し、これらを同じ解析で使うための設計上の工夫を説明する。

\subsection{機能群の役割と受け渡す量}
\begin{table}[htbp]
\centering\small
\begin{tabularx}{\linewidth}{p{24mm}YY}
\toprule
機能群 & 入力と主な処理 & 次へ渡す量・主な実装 \\
\midrule
設定・データ & case、種別データ、装置、観測の参照を展開し、単位と対応関係を検査 & 正規化した設定、状態・反応・放射記録。\code{io}, \code{data} \\
物理計算 & 領域別の電子状態と供給密度から率、準位収支、発光を評価 & 準位密度と体積分光発光。\code{physics}, \code{forward.emissivity} \\
視線・装置 & 発光を経路長で投影し、透過率、校正、LSF、画素平均を適用 & 視線・装置別の予測スペクトル。\code{geometry}, \code{instrument} \\
測定・推定 & 有効画素と波長窓を選び、共分散で残差を白色化して状態量を更新 & 名前付き推定量、予測・測定・残差。\code{inverse} \\
診断・評価 & 最終点の収支と数値品質、感度、profile、保留観測を評価 & 推定を支える情報とモデル改訂の根拠。\code{inverse}, \code{analysis} \\
\bottomrule
\end{tabularx}
\caption{計算機能の分担。\code{forward.compiled}が固定構造を準備し、\code{forward.model}が物理計算から観測予測までを接続する。各名称は\code{oescr}内の相対名である。}
\label{tab:function-groups}
\end{table}

設定・データの段階では、反応の参照先、密度を与える欄、装置IDと観測の対応を確定する。物理計算は、その設定を使って領域ごとの発光を返す。幾何と装置は、同じ発光を実験の視線・波長画素へ変換する。逆解析はこの予測値を測定と比較し、指定した状態量を更新する。この受け渡しを保つことで、例えば消光係数の改訂は準位密度と発光へ、LSFの改訂は測定予測へ、それぞれ対応する段階から反映される。

\subsection{順解析と逆解析の実行フロー}
図\ref{fig:computation-flow}に全体の流れを示す。順解析では、与えた状態に対して発光・観測の経路を一度評価する。逆解析では、その経路を目的関数として繰り返し呼び出し、状態量を測定に合わせる。

最初に、物理・装置の固定構造と、測定の画素状態・波長窓・共分散を準備する。次に、試行するパラメータから領域別のEEDFと密度を構成し、率積分、CR、発光、投影、装置応答を順に計算する。得られた予測と測定の残差を使って、最適化器が次の試行値を決める。終了解では物理量とスペクトルを標準形式で保存し、必要に応じて感度と数値品質を詳しく点検する。

\begin{figure}[p]
\centering\includegraphics[width=\linewidth]{\figfile{computation_flow.pdf}}
\caption{固定構造の準備、反復評価、最終出力を分けた計算フロー。左側は物理・観測の予測経路、右側は測定と推定の準備を表す。逆解析では、残差に基づくパラメータ更新が電子状態・密度の評価へ戻る。順解析は、この予測経路を指定した状態で一度評価する。最終診断は、推定結果を受けて実行する。$k$は反応速度係数、$\boldsymbol r$は白色化残差を表す。}
\label{fig:computation-flow}
\end{figure}

\subsection{固定構造の準備と反復時の再利用}
温度や電子密度の推定では、同じ状態集合、反応集合、視線、分光器を使って多数のスペクトルを計算する。そこで、\code{compile\_case}は入力の展開と参照検査をまとめ、\code{CompiledCase}に固定構造を保持する。状態の索引、反応の組立て規則、エネルギー・波長格子、原子線の基底、分子帯の輪郭、視線重み、観測格子への写像と画素積分の計画をこの段階で準備する。

反復時には、\code{plasma\_state}や\code{eedf}の値を更新し、EEDF、反応率、準位密度、発光振幅を再計算する。断面積を固定エネルギー格子へ補間した値は率計算器が保持する。反応、状態、幾何、装置等の構造を変更した場合は、対応する計算計画を再構築する。この扱いにより、係数やモデルを改訂した結果と、固定量の再利用とを整合させる。

\begin{table}[htbp]
\centering\small
\begin{tabularx}{\linewidth}{p{29mm}YY}
\toprule
処理の時期 & 主に準備・更新する量 & 数値計算上の役割 \\
\midrule
固定構造の準備 & 状態・反応の対応、格子、線・帯の形、視線と観測の写像 & 繰返しに共通する処理をまとめる \\
各試行状態の評価 & EEDF、外部密度、率、準位密度、発光、測定予測 & 状態量に依存する物理応答を更新する \\
最終点の評価 & 詳細収支、品質診断、感度、profile、局所曲率 & 得られた推定値の解釈と次の測定を支える \\
\bottomrule
\end{tabularx}
\caption{処理の頻度に応じた計算の分離。温度・密度の反復評価と、反応集合や装置を変更する再構成の役割を明確にする。}
\end{table}

\subsection{速度係数積分と準位系の数値解法}
エネルギー格子$E_q$上で、式\eqref{eq:rate}の積分核を$g_{r,q}=\sigma_r(E_q)v(E_q)\fe(E_q)$とすると、実装の台形則は
\begin{equation}
k_r^{(h)}=\sum_{q=0}^{N_E-2}\frac{E_{q+1}-E_q}{2}
\left(g_{r,q}+g_{r,q+1}\right)
\label{eq:discrete-rate}
\end{equation}
となる。断面積は固定格子への補間値を再利用し、電子速さとEEDFをその格子上で評価する。分布の更新によって積分核を変える。格子を細分化した比較は離散化への感度を、積分上限を広げた比較は高エネルギー成分への感度を調べる。これらを分けて評価すると、率の変化を格子間隔と積分範囲のどちらに帰属できるかを整理しやすい。

CRでは、各反応・放射・壁損失を一つずつ行列と右辺への寄与として組み立て、総和から式\eqref{eq:matrix}を作る。領域ごとに密行列の直接解法を適用するため、準位数$N_{\mathrm{lev}}$に対する主要な解法費用は概ね$N_{\mathrm{lev}}^3$に従う。選択準位を扱う縮約モデルに適した構成であり、追跡準位数は必要な励起・cascade経路に応じて決める。

同じ反応寄与から準位ごとの生成・損失収支を出力することも、本実装の特徴である。行列の係数、定常密度、反応ごとの収支を相互に照合でき、消光や段階励起を加えたときの作用を説明できる。数値品質の評価には、解の一意性、条件数、相対残差、非負化前の密度を用いる。

\subsection{発光格子、視線一括演算と画素積分}
固有発光では、原子線の解析的な規格化と、分子帯の元の支持域での規格化を準備する。領域ごとの密度と放射率で振幅を更新し、各成分を内部の細格子へ合成する。格子間隔は、固有線幅、帯の表間隔、LSF幅、画素間隔のうち、細かい表現が必要な尺度に合わせて選ぶ。

狭い原子線だけを扱うケースでは、観測装置の波長範囲と余白の和集合を使った分割格子を構成する。離れた観測帯域の間にある空白域の格子量を抑えられる。分子帯を含む場合や広い原子線の場合は、帯・線の支持域を保つ連続格子を使う。したがって、格子の選択は発光成分と観測帯域の両方に基づく。

領域発光を並べた配列へ視線行列$W$を掛けることで、複数視線を一括して投影する。同じ分光器を共有する視線には、波長軸を共通にした透過率・校正・LSFを適用する。LSFは各観測評価で核を作り、波長方向だけにFFT畳み込みを行う。この配列構成により、物理的な視線の違いと共通の装置応答を自然に扱える。

最後に、装置応答後の細格子値$y_q$を区分線形の関数として画素幅で積分する。$h_q=\lambda_{q+1}-\lambda_q$、$0\leq\xi\leq h_q$とすると、累積面積$H$は区間内で
\begin{equation}
H(\lambda_q+\xi)=H(\lambda_q)+\xi y_q+
\frac{\xi^2}{2h_q}\left(y_{q+1}-y_q\right)
\label{eq:antiderivative}
\end{equation}
と評価できる。各画素の境界$e_i,e_{i+1}$でこの累積面積を求め、
\begin{equation}
F_i=\frac{H(e_{i+1})-H(e_i)}{e_{i+1}-e_i}
\label{eq:bin-average-numeric}
\end{equation}
を出力する。画素境界が細格子点の間に来る場合にも、その部分区間の面積を含められる。未分解線や非等間隔画素を扱う際に、画素中心の点標本化に起因する積分誤差を抑え、線中心の移動に伴う各画素の光量変化を評価する処理である。

\subsection{共分散と補正量を共有した目的関数評価}
測定の準備では、有効画素、推定用の波長、観測と予測の対応を固定する。一般の測定共分散には、選択画素に対応するCholesky因子を保持して三角解法を用いる。装置gain、offset、tiltを推定する場合にも同じ画素と白色化を使い、物理量と補正量を一つの誤差基準へそろえる。

共通校正の乗算誤差には、式\eqref{eq:calcov}のrank-one構造を利用する。一つの校正groupについて、独立側の共分散を$C_0=LL^{\mathsf T}$、その白色化を$B=L^{-1}$、校正ベクトルを$\boldsymbol u$とする。$\boldsymbol w=B(\boldsymbol F-\boldsymbol y)$、$\boldsymbol v=B\boldsymbol u$、$a=\sqrt{1+\boldsymbol v^{\mathsf T}\boldsymbol v}$とおけば、補正後残差は
\begin{equation}
\boldsymbol r=\boldsymbol w-
\frac{\boldsymbol v\left(\boldsymbol v^{\mathsf T}\boldsymbol w\right)}{a(a+1)}
\label{eq:rank-one-whitening}
\end{equation}
と書ける。これは$(I+\boldsymbol v\boldsymbol v^{\mathsf T})^{-1/2}\boldsymbol w$の評価に当たる。共通の尺度誤差が繰返し測定へ及ぼす相関を保ちつつ、その追加のために全画素の密な共分散を再分解する費用を抑える。

比や帯面積も使う場合は、選択したfeatureとその共分散を目的関数へ組み込む。スペクトルだけで補正量が線形なら重み付き線形最小二乗を、featureや事前項を含む場合には対応する内側最適化を使う。この内側評価を含めた残差を外側へ返すことで、温度・密度に対する感度も実際に使う目的関数から求められる。

\subsection{最適化と最終診断の分離}
密度等の正の量は$\log_{10}$座標を選べるため、桁の異なる量を有限範囲へ写して推定できる。広い範囲を探索するときはdifferential evolutionで候補点を求め、その後に境界付き最小二乗で精密化する。局所推定だけを使う構成も選べる。計算時間と探索範囲は、推定量数と初期値についての知識に応じて設定する。

反復中は予測スペクトルと残差の計算を中心に実行する。実行時に\code{--audit}を指定すると、終了解で順解析品質、測定だけのJacobian、特異値、条件付きprofile、局所曲率を求める。測定残差数がパラメータ数以上のJacobianにはeconomy SVDを使い、観測空間の大きな正方行列の構成を抑える。パラメータ数が測定残差数を上回る場合は、右零空間の方向も保持する。

この分離は、反復計算の作業量を抑えるだけでなく、結果を読む順序も明確にする。まず測定との一致を確認し、次に観測が支えるパラメータの組合せを調べ、その後に事前分布を含む局所不確かさを評価する。複数の線・視線を追加する効果も、この同じ順序で比較できる。

\subsection{本実装の工夫と検証との対応}
\begin{table}[htbp]
\centering\small
\begin{tabularx}{\linewidth}{p{34mm}YY}
\toprule
設計上の工夫 & 対象となる数値上の課題 & 確認方法 \\
\midrule
固定構造と状態更新の分離 & 繰返しの準備費用、構造変更との整合 & 計画再利用と構造変更時の再構築を確認 \\
発光基準の規格化と細格子 & 観測区間の変更、未分解線の積分 & 解析光量、Gaussian線と画素位置の比較 \\
過程寄与からの行列・収支生成 & 供給・損失の符号と表示の整合 & 二準位・cascadeの解析解、寄与の総和 \\
共通白色化と校正相関の利用 & 物理量とgainの重みの整合、共通尺度誤差 & 直接構成した共分散による評価と比較 \\
推定と測定感度の段階的評価 & 状態量の縮退、事前条件の影響の把握 & 既知の縮退問題、生成回復、profile \\
\bottomrule
\end{tabularx}
\caption{計算の見通しと一貫性を保つための実装上の選択。解析解、保存則、共分散の一致、既知の感度問題を使って、それぞれの性質を確認する。}
\label{tab:numerical-design}
\end{table}

これらの工夫は、物理データを変更した際の影響を反応収支から測定残差まで追跡し、同じ観測基準で推定結果を比較するために組み合わせている。第\ref{sec:benchmarks}節では、この計算経路を具体的な問題設定に適用し、積分量と状態量の回復を評価する。
